\chapter{Теоретические основы прогнозирования финансовых временных рядов}

\section{Особенности финансовых временных рядов}
Временным рядом (или рядом динамики) называется последовательность значений некоторого статистического показателя (случайной величины), упорядоченных в хронологическом порядке — в последовательные моменты или периоды времени. В экономике и финансах типичными примерами временных рядов выступают курсы обмена валют, котировки акций, уровень инфляции, показатели спроса и валовой национальный продукт. В зависимости от характера фиксации времени ряды могут быть дискретными и непрерывными. В непрерывных рядах значения переменной фиксируются непрерывно во времени, однако на практике (особенно в финансах) чаще всего анализируются дискретные временные ряды, где наблюдения производятся через определенные фиксированные интервалы времени (например, ежечасно, ежедневно, ежемесячно). Таким образом, временной ряд в контексте данной работы — это упорядоченная во времени последовательность исторических значений валютного курса.

В классическом анализе любой экономический временной ряд традиционно рассматривается как результат взаимодействия нескольких основных компонент:
\begin{itemize}
    \item Тренд (тенденция) — плавно меняющаяся компонента, описывающая чистое влияние долговременных факторов и отражающая общую (вековую) тенденцию процесса к росту или снижению.
    \item Сезонность — регулярные колебания, отражающие внутригодовую или внутримесячную повторяемость экономических процессов.
    \item Цикличность — последовательность волнообразных флуктуаций (экономических циклов), длительность которых обычно превышает один год (например, циклы деловой активности).
    \item Случайная (нерегулярная) составляющая (шум) — непредсказуемые колебания, отражающие влияние множества не поддающихся учету случайных факторов.
\end{itemize}


Ключевым свойством, необходимым для анализа и моделирования, является стационарность временного ряда. Строго стационарным называется ряд, совместное распределение вероятностей которого не изменяется при сдвиге во времени. В практических исследованиях чаще опираются на понятие слабой стационарности: ряд считается стационарным, если его математическое ожидание (среднее значение) и дисперсия остаются постоянными во времени, а ковариация зависит только от временного лага.

Стационарность крайне важна, поскольку большинство классических статистических методов и теорем опираются на предположение о неизменности свойств процесса. Проблема заключается в том, что реальные финансовые ряды практически всегда нестационарны. Их среднее значение и дисперсия меняются с течением времени, формируя стохастические тренды. Для приведения таких рядов к стационарному виду обычно применяются математические преобразования, чаще всего — переход от абсолютных значений к их разностям (приращениям).

Финансовые временные ряды, в частности курсы валют, обладают ярко выраженной спецификой, которая значительно усложняет их математическое описание:
\begin{itemize}
    \item Волатильность. Для финансовых рынков характерны резкие скачки и нестабильность дисперсии. Широко известен эффект «кластеризации волатильности»: периоды сильных и резких колебаний цены часто сменяются периодами относительного затишья.
    \item Нестационарность. Валютные курсы подвержены влиянию макроэкономических трендов и структурных изменений. Часто их поведение близко к процессу «случайного блуждания» (random walk), что делает их непредсказуемыми на больших горизонтах.
    \item Автокорреляция. Текущие значения финансовых показателей и их волатильность часто зависят от прошлых значений (наличие "памяти" у рынка), что фиксируется с помощью автокорреляционных функций.
    \item Нелинейность. Динамика финансовых рынков характеризуется высокой сложностью, а поведение участников рынка приводит к тому, что изменения цен не описываются простыми линейными формулами.
    \item Шум и случайность. Финансовые показатели содержат значительную долю шума (соотношение сигнал/шум крайне низкое). На них ежеминутно влияют поступление новых случайных новостей, макроэкономическая статистика и политические события.
\end{itemize} 

Специфические свойства финансовых временных рядов порождают серьезные трудности при их прогнозировании. Из-за беспрецедентно высокой неопределенности, динамизма и хаотичности финансовой среды предсказание курсов валют является задачей повышенной сложности. Традиционные статистические подходы и линейные модели (такие как ARIMA или экспоненциальное сглаживание), несмотря на их полезность, не всегда справляются с нелинейной природой и высокой зашумленностью финансовых данных. Это обуславливает необходимость использования более сложных и гибких аналитических инструментов.

Финансовые временные ряды представляют собой сложные, хаотичные и нестационарные структуры. Присущие им специфические свойства — высокая волатильность, нелинейность, подверженность случайным шокам и структурным сдвигам — требуют для анализа и точного прогнозирования применения современных продвинутых методов, в том числе алгоритмов машинного обучения и нейронных сетей.

\section{Классические модели прогнозирования курса валют}
% TODO

\section{Методы машинного обучения в прогнозировании}
% TODO

\section{Нейросетевые модели прогнозирования}
% TODO
